% BEGIN VOL08-EXPANSION VIII25
\section{Supplementary worked examples}
\label{sec:viii25-pedagogy}

\begin{example}[The circle with arbitrary coefficients]\label{ex:viii25-ped-01}
For any abelian group $G$,
\[
H_0(S^1;G)\cong G,\qquad H_1(S^1;G)\cong G,
\]
and higher groups vanish. The cellular chain complex has one copy of $G$ in degrees $0$ and $1$ with zero differential.
\end{example}

\begin{example}[Torsion changes under coefficients]\label{ex:viii25-ped-02}
For $\mathbb{RP}^2$, the cellular differential $C_2\to C_1$ is multiplication by $2$. With $\mathbb F_2$ coefficients it becomes zero, so mod-$2$ homology sees a degree-$2$ class that integral homology does not.
\end{example}

\begin{example}[Tensor and Tor effects]\label{ex:viii25-ped-03}
Changing coefficients is not merely tensoring homology in all cases. The universal coefficient theorem gives
\[
0\to H_n(X;\mathbb Z)\otimes G\to H_n(X;G)\to
\operatorname{Tor}(H_{n-1}(X;\mathbb Z),G)\to0,
\]
so torsion one degree lower can create new classes.
\end{example}

\section{Graded supplementary exercises}
The following exercises progress from direct computation and construction to proof, hypothesis testing, application, and synthesis.

\subsection*{Standard computations and constructions}

\begin{exercise}[Point coefficients]\label{exr:viii25-ped-01}
Compute $H_0(\{*\};G)$.
\end{exercise}
\begin{hint}
The degree-zero chain group is $G$.
\end{hint}
\begin{solution}
$H_0(\{*\};G)\cong G$.
\end{solution}

\begin{exercise}[Circle coefficients]\label{exr:viii25-ped-02}
Compute $H_1(S^1;G)$.
\end{exercise}
\begin{hint}
Use the one-cell CW complex.
\end{hint}
\begin{solution}
$H_1(S^1;G)\cong G$.
\end{solution}

\begin{exercise}[Sphere coefficients]\label{exr:viii25-ped-03}
Compute $H_n(S^n;G)$ for $n>0$.
\end{exercise}
\begin{hint}
Use one cell in degrees $0$ and $n$.
\end{hint}
\begin{solution}
$H_n(S^n;G)\cong G$, with $H_0\cong G$ and no other nonzero groups.
\end{solution}

\begin{exercise}[Tensor with Z/m]\label{exr:viii25-ped-04}
Compute $\mathbb Z/n\otimes\mathbb Z/m$.
\end{exercise}
\begin{hint}
The answer has order $\gcd(m,n)$.
\end{hint}
\begin{solution}
$\mathbb Z/n\otimes\mathbb Z/m\cong\mathbb Z/\gcd(m,n)$.
\end{solution}

\begin{exercise}[Tor with Z/m]\label{exr:viii25-ped-05}
Compute $\operatorname{Tor}(\mathbb Z/n,\mathbb Z/m)$.
\end{exercise}
\begin{hint}
Use a two-term free resolution of $\mathbb Z/n$.
\end{hint}
\begin{solution}
It is also $\mathbb Z/\gcd(m,n)$.
\end{solution}

\subsection*{Proofs}

\begin{exercise}[Chain-level definition]\label{exr:viii25-ped-06}
Explain why $C_*(X;G)=C_*(X;\mathbb Z)\otimes G$ is again a chain complex.
\end{exercise}
\begin{hint}
Tensor the boundary with the identity.
\end{hint}
\begin{solution}
The differential is $\partial\otimes\operatorname{id}_G$, whose square is $(\partial^2)\otimes\operatorname{id}=0$.
\end{solution}

\begin{exercise}[Field simplification]\label{exr:viii25-ped-07}
Show that over a field, homology can be computed by ordinary linear algebra.
\end{exercise}
\begin{hint}
Chain groups become vector spaces.
\end{hint}
\begin{solution}
Boundary maps are matrices over the field, so $H_n=\ker d_n/\operatorname{im}d_{n+1}$ is obtained from ranks and nullspaces.
\end{solution}

\begin{exercise}[Naturality in coefficients]\label{exr:viii25-ped-08}
Explain how a homomorphism $G\to G'$ induces $H_n(X;G)\to H_n(X;G')$.
\end{exercise}
\begin{hint}
Tensor the chain complex with the coefficient homomorphism.
\end{hint}
\begin{solution}
The map $\operatorname{id}\otimes\phi$ is a chain map and therefore induces a homology map.
\end{solution}

\begin{exercise}[Mod-p detection]\label{exr:viii25-ped-09}
Why can mod-$p$ coefficients reveal $p$-torsion?
\end{exercise}
\begin{hint}
Both tensor and Tor terms can survive.
\end{hint}
\begin{solution}
$p$-torsion contributes to tensor products in its own degree and through Tor from the preceding degree, producing nonzero mod-$p$ classes.
\end{solution}

\subsection*{Counterexamples and hypothesis tests}

\begin{exercise}[Coefficients never change Betti numbers]\label{exr:viii25-ped-10}
Test: changing coefficients never changes homology dimensions.
\end{exercise}
\begin{hint}
Compare integral and mod-$2$ homology of $\mathbb{RP}^2$.
\end{hint}
\begin{solution}
False. Torsion can become visible over finite fields and change dimensions.
\end{solution}

\begin{exercise}[Tensor only]\label{exr:viii25-ped-11}
Test: $H_n(X;G)$ is always $H_n(X;\mathbb Z)\otimes G$.
\end{exercise}
\begin{hint}
Remember the Tor correction.
\end{hint}
\begin{solution}
False; the universal coefficient theorem includes $\operatorname{Tor}(H_{n-1}(X),G)$.
\end{solution}

\begin{exercise}[Field sees all torsion orders]\label{exr:viii25-ped-12}
Test: homology over one field determines all integral torsion.
\end{exercise}
\begin{hint}
A field only probes its characteristic.
\end{hint}
\begin{solution}
False. For example, $\mathbb F_p$ is sensitive to $p$-primary behavior but not to torsion of unrelated prime order.
\end{solution}

\subsection*{Applications and investigations}

\begin{exercise}[Choosing coefficients]\label{exr:viii25-ped-13}
Why might one compute first over several finite fields before doing Smith normal form over $\mathbb Z$?
\end{exercise}
\begin{hint}
Field linear algebra is cheaper.
\end{hint}
\begin{solution}
Finite-field computations quickly reveal Betti numbers and likely torsion primes, guiding or validating the more expensive integral computation.
\end{solution}

\begin{exercise}[Orientation coefficients]\label{exr:viii25-ped-14}
Why do coefficient choices matter for orientability questions?
\end{exercise}
\begin{hint}
A sign obstruction disappears in characteristic $2$.
\end{hint}
\begin{solution}
Over $\mathbb F_2$, $-1=1$, so every manifold has a mod-$2$ fundamental class even when no integral orientation exists.
\end{solution}

\subsection*{Challenge problems}

\begin{exercise}[RP2 mod 2]\label{exr:viii25-ped-15}
Using its cellular differential $2$, compute $H_*(\mathbb{RP}^2;\mathbb F_2)$.
\end{exercise}
\begin{hint}
Multiplication by $2$ becomes zero.
\end{hint}
\begin{solution}
There is one $\mathbb F_2$ in degrees $0,1,2$.
\end{solution}

\begin{exercise}[Coefficient synthesis]\label{exr:viii25-ped-16}
Describe how integral, rational, and mod-$p$ homology complement one another.
\end{exercise}
\begin{hint}
Separate free rank from prime torsion information.
\end{hint}
\begin{solution}
Rational homology isolates free ranks; mod-$p$ homology cheaply detects $p$-primary effects; integral homology assembles the full free and torsion structure.
\end{solution}

% END VOL08-EXPANSION VIII25
